\worldchapter{\trans{recipes_title}}{brewing}
\label{ch:appendix-recipes}

\begin{multicols}{2}


    \section{Tränke und Tinkturen}



    \subsection*{Lachtrank}

    \begin{tabular}{@{}ll@{}}
        \textbf{\trans{difficulty_label}:} & Einfach \\
        \textbf{\trans{yield_label}:} & 1 \\
    \end{tabular}

    Zubereitung:
Zuerst müsst ihr das Quellwasser in einem Kupferkessel zum Sieden bringen, bis es sanft vor sich hin blubbert. Gebt den zerstoßenen Beinwell hinzu, er bildet die physische Basis und sorgt dafür, dass die Lachmuskeln geschmeidig bleiben.

Nun kommt der entscheidende Moment: Die „Kicher-Erbse“ muss unter lautem Erzählen eines (wirklich!) schlechten Witzes im Mörser zerstoßen werden. Je flacher der Witz, desto spritziger der Trank! Rührt das Ganze dreimal gegen den Uhrzeigersinn, während ihr dabei leise kichert. Zum Schluss wird der Honig eingerührt, um die Bitterkeit der Welt zu vertreiben.

Zubereitungszeit: 30 Minuten aktives Rühren und Witzeerzählen, danach 1 Stunde Ruhezeit zum „Abklingeln“.

Menge: 1 Anwendung, ca. 50 Oth Flüssigkeit, Gewicht: 50 Oth.

Wirkung:
Der Konsument muss bei nächster Gelegenheit über jeden noch so kleinen Vorfall lauthals lachen (Dauer: 2W6 Minuten). Spielmechanisch verleiht dies einen Bonuswürfel auf alle Proben zur „Darbietung“ oder „Kommunikation“, erschwert jedoch alle „Heimlichkeits“-Proben um +3 (SG Mittel).

Erschöpfung: Nach Ablauf der Wirkung -1 Punkt auf alle Würfelproben für 30 Minuten (Lachen ist anstrengend!).

    \paragraph{\trans{ingredients_label}}
    \begin{denseitemize}
        \item \textbf{Wasser} (60 Oth), \textit{Zutaten}.
        Kaltes klares Wasser.
        \item \textbf{Beinwell (Symphytum officinale)} (10 Oth), \textit{Zutaten}.
        Beinwell regt die Durchblutung an, Prellungen, Hämatome und Verstauchungen verschwinden schneller. Beinwell beschleunigt die Regeneration der Zellen.
        \item \textbf{Kichererbsen} (1 Stück), \textit{Nahrung / Proviant}.
        Diese kleine, gelb-grüne Hülsenfrucht ist ein wahrlich wundervolles Nahrungsmittel.

Abgesehen von ihrer alchemistischen Verwendung als Basis für Lachtränke, ist sie eine hervorragende Speise für müde Wanderer.

Man findet sie oft auf den reichhaltigen Märkten in und um Al Bah Ji Ra. Ihre Beliebtheit verbreitet sich jedoch langsam in ganz Tirakan.

Der Verzehr einer Handvoll gerösteter Kicher-Erbsen vertreibt trübe Gedanken und lässt den Schatten im Gemüt für kurze Zeit schrumpfen. Vor allem wenn man sie mit Gewürzen wie Knoblauch, Kümmel oder Pfeffer würzt.

Auch pur sind sie reich an Proteinen, komplexen Ballaststoffen und dienen so als nahrhafte Ergänzung für jeden Speiseplan.

„Wenn der Eintopf zu bitter schmeckt und das Leben zu schwer wiegt, wirf eine Handvoll Kichererbsen hinein. Dein Magen wird es dir mit einem Glucksen danken.“
— Meisterkoch Alar al-Din des Zelts der Sieben Schleier von El Kurru
        \item \textbf{Honig} (5 Oth), \textit{Nahrung / Proviant}.
        In den kultivierten Regionen wird der Honig von fleißigen Imkern gewonnen, die ihre Stöcke oft in der Nähe von Kräutergärten aufstellen.

Dieser Honig ist klar, beständig und hat einen feinen Geschmack nach Thymian oder Lavendel. Er kostet in der Regel etwa 2 Gulden pro 5 Oth und ist die verlässliche Basis für jeden Alchemisten, der die Bitterkeit von Kräuterextrakten mildern will.

Wer es wagemutiger mag, sucht in den Wäldern nach den Nestern der Wildbienen. Dieser Honig ist dunkler, zäher und oft mit Pollen oder kleinen Wachsstücken versetzt. Er hat eine kräftige, fast erdige Note. Man sagt, wilder Honig aus den Urwäldern besitze eine stärkere regenerative Kraft für die Stimme. Ideal für Barden, die nach einer langen Nacht in der Schenke ihre Kehle ölen müssen.
    \end{denseitemize}


    \subsection*{Kinstarchel Sekret}

    \begin{tabular}{@{}ll@{}}
        \textbf{\trans{difficulty_label}:} & Mittel \\
        \textbf{\trans{yield_label}:} & 1 \\
    \end{tabular}

    Zubereitung:
Zuerst müssen die Knochen eines frisch (oder zumindest nicht allzu verwitterten) verstorbenen Kinstarchels sorgsam gereinigt und zerschlagen werden, um an das tiefschwarze Mark zu gelangen. Dieses wird in reinem Alkohol, über kleiner Flamme erhitzt. Man darf es nicht sieden lassen! Sobald sich eine ölige, leicht schimmernde Schicht an der Oberfläche bildet, muss diese mit einer silbernen Pipette abgezogen werden. Das Sekret wird dann direkt mit einem bereits fertigen Trank oder Gift vermischt. Die Mischung ist hochgradig instabil; eine heftige Erschütterung (wie der Aufprall einer Phiole) genügt, um die enthaltene Energie schlagartig freizusetzen.

Zubereitungszeit: 1,5 Stunden aktive Extraktion, danach muss die Mischung 1 Stunde in einem kühlen, dunklen Gefäß ruhen, bevor sie stabil genug für den Transport ist.

Menge: 1 Anwendung (ca. 10 Oth Sekret), ausreichend, um eine Phiole (50-100 Oth) eines anderen Tranks zu präparieren.

Wirkung:
Das Sekret verleiht einem Trank die Eigenschaft „Explosiv“. Beim Zerbrechen des Gefäßes (Wurf-Probe auf Gewandtheit erforderlich) explodiert die Mischung in einem Umkreis von 1W3 Metern. Jedes Wesen in diesem Bereich wird so behandelt, als hätte es den Trank oder das Gift direkt eingenommen.

Anmerkung:
Das Sekret selbst verursacht keinen Schaden, es dient lediglich als rasanter Trägerstoff für die Wirkung des beigemischten Tranks.

    \paragraph{\trans{ingredients_label}}
    \begin{denseitemize}
        \item \textbf{Alkohol} (300 Oth), \textit{Zutaten}.
        In den Laboren der Alchemisten und den Hütten der Kräuterkundigen wird Alkohol selten zum Vergnügen gelagert. Er gilt als Lösungsmittel, dazu fähig, die Essenz von Pflanzen und Mineralien zu entfesseln.

Er wird durch mehrfache Destillation aus vergorenem Getreide oder Früchten gewonnen. Esist ein brennend scharf riechendes Destillat, das so stark ist, dass es beim Einatmen die Nase kitzelt und auf der Haut sofort kühl verdunstet.
        \item \textbf{Kinstarchel Knochen} (2 Stein), \textit{Zutaten}.
        Knochen eines verstorbenen Kinstarchels. Aus dem Mark der Knochen kann ein explosives Sekret gewonnen werden.
    \end{denseitemize}


    \subsection*{Trank der tiefen Ruhe}

    \begin{tabular}{@{}ll@{}}
        \textbf{\trans{difficulty_label}:} & Mittel \\
        \textbf{\trans{yield_label}:} & 2 \\
    \end{tabular}

    Zubereitung:
Das Brauen dieses Tranks beginnt damit, dass das Wasser zum Kochen gebracht wird, um anschließend Baldrian und zerstoßene Monddornbeeren für fünf Minuten darin zu kochen. Der Sud wird danach vom Feuer genommen und muss dreißig Minuten ziehen, bevor er abgefüllt wird.

Zubereitungszeit: 50 Minuten, 6 Stunden Ruhezeit

Menge: 2 Anwendungen, 2x250 Oth Flüssigkeit, 2x250 Oth Gewicht

Wirkung:
Die Essenz führt dazu, dass der Charakter innerhalb von fünf Minuten in einen tiefen, zehnstündigen Schlaf fällt, wobei währenddessen W6 Wunden regeneriert werden. Um aus diesem tiefen Schlummer zu erwachen, muss der Charakter bei Störungen oder Weckversuchen zwei erfolgreiche Würfe auf Resistenz in Folge erzielen. Nur direkte körperliche Gewalt (Schütteln/Angriff) oder extrem laute Geräusche bewirken ein sofortiges Erwachen.

    \paragraph{\trans{ingredients_label}}
    \begin{denseitemize}
        \item \textbf{Monddornbeere} (5 Beeren), \textit{Zutaten}.
        Die Monddornbeere ist ein Geschenk der tiefsten Nacht. Sie wächst als bodendeckender Strauch, dessen zarte Ranken und tiefgrüne Blätter von markanten, kurzen Dornen geschützt werden.

Man findet sie ausschließlich an Orten, die selten das Licht der Sonne erblicken, meist tief in uralten Wäldern oder in der Nähe von feuchten Grotten- und Höhleneingängen. Ihre wahre Pracht entfaltet sie nur unter dem Licht des Vollmonds, wenn ihre kleinen, beerenartigen Früchte in einem geheimnisvollen, matten Blau leuchten, fast als hätten sie das Licht der Himmelskugel selbst verschluckt.

Die Beere ist berüchtigt für ihre starke sedierende Wirkung. In kleiner Dosis wirkt sie beruhigend, doch konzentriert in einem Trank, kann ihre Essenz den Geist betäuben und den Körper in einen Zustand tiefer, traumloser Stille versetzen. Das Sammeln ist riskant, da ihre Dornen bei Berührung einen temporären Juckreiz auslösen können.
        \item \textbf{Baldrian (Valeriana officinalis)} (1 Bund), \textit{Zutaten}.
        Baldrian hilft bei Einschlafstörungen und Unruhe. Hopfen und Melisse steigern die Wirkung des Baldrians und verbessern den Geschmack.
        \item \textbf{Wasser} (500 Oth), \textit{Zutaten}.
        Kaltes klares Wasser.
    \end{denseitemize}


    \subsection*{Einfache Wundtinktur}

    \begin{tabular}{@{}ll@{}}
        \textbf{\trans{difficulty_label}:} & Einfach \\
        \textbf{\trans{yield_label}:} & 1 \\
    \end{tabular}

    Zubereitung:Zuerst müsst Ihr den Alant und den Beinwell sorgfältig reinigen. Die Wurzeln des Alants werden in feine Scheiben geschnitten und in einem Mörser zu einem Brei zerstoßen, während die Blätter des Beinwells nur leicht zwischen den Fingern gerieben werden, um ihre Säfte freizusetzen. Gebt beides in eine Keramikflasche und füllt sie mit klarem Quellwasser (oder, wenn Ihr es besonders rein wollt, mit destilliertem Geist) auf. Lasst das Gemisch über einem sanften Feuer köcheln, bis die Flüssigkeit eine tiefgrüne, fast ölige Konsistenz annimmt. Seid vorsichtig! Wer das Feuer zu heiß schürt, verbrennt die heilenden Geister der Pflanzen!

Zubereitungszeit:
Aktive Zeit: 30 Stunden (Zerkleinern und Sieden).Ruhezeit: 2 Stunden an einem kühlen, dunklen Ort, damit die Wirkstoffe in die Tinktur übergehen können.

Menge:
1 Anwendung, ca. 10 Oth Flüssigkeit in einer Flasche.
Gesamtgewicht: 15 Oth (inklusive Gefäß).

Wirkung:
Bei erfolgreicher Anwendung mit Erster Hilfe und einer Bandage heilt die Bandage 1W3 Wunden zusätzlich.

    \paragraph{\trans{ingredients_label}}
    \begin{denseitemize}
        \item \textbf{Wasser} (100 Oth), \textit{Zutaten}.
        Kaltes klares Wasser.
        \item \textbf{Alant (Inula helenium)} (10 Oth), \textit{Zutaten}.
        Diese Heilpflanze aus dem Mittelalter ist in modernen Zeiten nicht mehr weit verbreitet. Ihre Anwendung verbessert die Verdauung, und ihr wird eine vorbeugende Wirkung gegen Darmkrebs zugeschrieben.
        \item \textbf{Beinwell (Symphytum officinale)} (10 Oth), \textit{Zutaten}.
        Beinwell regt die Durchblutung an, Prellungen, Hämatome und Verstauchungen verschwinden schneller. Beinwell beschleunigt die Regeneration der Zellen.
    \end{denseitemize}


    \subsection*{Zaubertrank (eine volle Karaffe)}

    \begin{tabular}{@{}ll@{}}
        \textbf{\trans{difficulty_label}:} & Einfach \\
        \textbf{\trans{yield_label}:} & 3 \\
    \end{tabular}

    Zubereitung:
Zuerst bereitet Ihr den Sud aus Alant, Beinwell und Schafgarbe vor. Wenn der Dampf beginnt, schwer und kräuterig in der Nase zu brennen, fügt Ihr die Schuppe einer Flussjungfer hinzu. Rührt siebenmal gegen den Sonnenlauf, während Ihr an das kalte, tiefe Wasser denkt. Der Trank wird nicht leuchten wie ein billiger Jahrmarktstrick, sondern ein tiefes, erdiges Blau annehmen, das im richtigen Winkel schimmert.

Zubereitungszeit:
Aktive Zeit: 30 Minuten (das Schaben und Reinigen der Schuppe dauert seine Zeit)
Ruhezeit: 2 Stunden 

Menge:
5 Anwendungen, 500 Oth Flüssigkeit, Gewicht: 500 Oth

Wirkung: Stellt pro Anwendung 2 Arkana wieder her.

    \paragraph{\trans{ingredients_label}}
    \begin{denseitemize}
        \item \textbf{Alant (Inula helenium)} (10 Oth), \textit{Zutaten}.
        Diese Heilpflanze aus dem Mittelalter ist in modernen Zeiten nicht mehr weit verbreitet. Ihre Anwendung verbessert die Verdauung, und ihr wird eine vorbeugende Wirkung gegen Darmkrebs zugeschrieben.
        \item \textbf{Beinwell (Symphytum officinale)} (15 Oth), \textit{Zutaten}.
        Beinwell regt die Durchblutung an, Prellungen, Hämatome und Verstauchungen verschwinden schneller. Beinwell beschleunigt die Regeneration der Zellen.
        \item \textbf{Schafgarbe (Achillea millefolium)} (5 Oth), \textit{Zutaten}.
        Schafgarbe wird wegen ihrer blutstillenden Wirkung eingesetzt. Die Blüten und die Blätter enthalten Gerb-, Bitter- und Mineralstoffe. Das ätherische Öl der Pflanze wirkt entzündungshemmend und krampflösend.
        \item \textbf{Schuppe einer Flussjungfer} (1 Stück), \textit{Zutaten}.
        Die Flussjungfer ist eine überaus scheue Kreatur, die, so munkelt man in Alchemistischen Kreisen, wohl noch kein sterbliches Wesen in Tirakan wahrhaftig zu Gesicht bekommen hat. Der einzige greifbare Beweis ihrer Existenz sind die Schuppen.

Man findet sie gelegentlich an schlammigen Flussufern, auf gischtumspülten Felsen oder in der Finsternis unterirdischer Seen entdeckt.

In der Dämmerung wirken diese Schuppen oft schlicht grün-bräunlich, fast wie herkömmliches Horn oder vertrocknetes Laub. Doch sobald die Sonne Tirakans die Oberfläche küsst, erwachen sie zum Leben und schimmern in allen Farben des Perlmutt. Sie erreichen oft die Größe einer stolzen Handfläche und ähneln in ihrer Beschaffenheit den Schuppen eines großen Fisches.

Glücklicherweise ist diese magische Zutat auf den Märkten des Reichs, besonders auf den Handelsrouten zwischen Toran und Yavon bis hinunter nach Meridian, kein seltener Anblick. Da sie regelmäßig gefunden werden, bleibt ihr Wert überschaubar, was sie zu einer ehrlichen Zutat macht. Sie sind sogar so verbreitet, dass es sich für Fälscher kaum lohnt, minderwertige Kopien zu fertigen. Aber ein Wort der Warnung von mir: Achtet stets auf das charakteristische Schimmern! 

Einer alten Legende nach thronen die Flussjungfern in einsamen Nächten auf bemoosten Felsen und richten ihr endloses, wasserfarbenes Haar mit Kämmen aus Knochen und altem Treibholz. Ihr Antlitz sei dabei von einer friedlichen Melancholie gezeichnet, während sie leise Lieder summen, deren tiefere Bedeutung allein dem fließenden Wasser vorbehalten bleibt.
        \item \textbf{Wasser} (500 Oth), \textit{Zutaten}.
        Kaltes klares Wasser.
    \end{denseitemize}


    \subsection*{Gilden Elexir}

    \begin{tabular}{@{}ll@{}}
        \textbf{\trans{difficulty_label}:} & Einfach \\
        \textbf{\trans{yield_label}:} & 3 \\
    \end{tabular}

    Zubereitung:
Zuerst wird das reine Quellwasser erhitzt, nicht kochen, nur sieden! Das Wachkraut wird fein gehackt und zusammen it dem Honig in as heiße Wasser gegeben.

Die Mischung muss so lange ziehen, bis der Honig sich vollständig aufgelöst hat und das Wasser die goldene Farbe
der Kräuter annimmt. Danach wird das Ganze durch ein feines Leinentuch in die Phiole abgegossen.

Zubereitungszeit:
1 Stunde aktives Brauen am Kessel,
gefolgt von einer halben Stunde Auskühlzeit.

Menge:
3 Anwendungen. 200 Oth Flüssigkeit, Gesamtgewicht 250 Oth (inklusive Glasphiole).

Wirkung:
Bei Einnahme einer Anwenung wird sofort vorhandene Erschöpfung abgebaut. Zudem erhält der Konsument für die nächsten 1W6 Stunden einen Bonus von +1 Würfel auf alle Proben, die Konzentration oder Wachsamkeit er-
fordern.

    \paragraph{\trans{ingredients_label}}
    \begin{denseitemize}
        \item \textbf{Honig} (20 Oth), \textit{Nahrung / Proviant}.
        In den kultivierten Regionen wird der Honig von fleißigen Imkern gewonnen, die ihre Stöcke oft in der Nähe von Kräutergärten aufstellen.

Dieser Honig ist klar, beständig und hat einen feinen Geschmack nach Thymian oder Lavendel. Er kostet in der Regel etwa 2 Gulden pro 5 Oth und ist die verlässliche Basis für jeden Alchemisten, der die Bitterkeit von Kräuterextrakten mildern will.

Wer es wagemutiger mag, sucht in den Wäldern nach den Nestern der Wildbienen. Dieser Honig ist dunkler, zäher und oft mit Pollen oder kleinen Wachsstücken versetzt. Er hat eine kräftige, fast erdige Note. Man sagt, wilder Honig aus den Urwäldern besitze eine stärkere regenerative Kraft für die Stimme. Ideal für Barden, die nach einer langen Nacht in der Schenke ihre Kehle ölen müssen.
        \item \textbf{Wasser} (130 Oth), \textit{Zutaten}.
        Kaltes klares Wasser.
        \item \textbf{Gildenkraut (Vigilia mercatoris)} (1 Bund), \textit{Zutaten}.
        Das Gildenkraut ist eine zähe, beinahe unscheinbare Pflanze, die nur dem geschulten Auge eines Kundigen auffällt.

Sie zeichnet sich durch gezackte, ledrige Blätter von tiefgrüner Färbung aus, deren dicke Blattadern einen leichten rötlichen Schimmer aufweisen. Ihre Blüten sind winzig, blassgelb und schließen sich sofort, sobald die Sonne ihren Zenit überschreitet.

Das Kraut ist vollkommen mundan und besitzt nicht den geringsten Funken Arkana. Sein Wert liegt vielmehr in der extrem hohen, natürlichen Konzentration an Alkaloiden, vor allem reinem Koffein. Zerreibt man die harten Blätter zwischen den Fingern, verströmen sie einen scharfen, erdigen Duft, der an stark geröstete Nüsse erinnert.

Lage und Lebensraum:
Dieses Gewächs ist verhältnismäßig selten, da es ganz bestimmte, widrige Umstände benötigt, um seine anregenden Säfte zu bilden. Es gedeiht ausschließlich in kargen, mineralreichen Böden und meidet das feuchte Moos tiefer Wälder.

Man findet das Gildenkraut vornehmlich in den felsigen Spalten uralter Handelsrouten, im bröckelnden Steinmörtel verlassener Gilden-Außenposten oder an den steilen, sonnengegerbten Hängen des Toraner Hochlands. Es braucht den harten Fels und viel Wind, um nicht zu verkümmern.
    \end{denseitemize}


    \subsection*{Starker Heiltrank}

    \begin{tabular}{@{}ll@{}}
        \textbf{\trans{difficulty_label}:} & Schwer \\
        \textbf{\trans{yield_label}:} & 3 \\
    \end{tabular}

    Zubereitung:
Zunächst muss das reine Quellwasser über einem Feuer aus getrocknetem Eichenholz langsam zum Sieden gebracht werden. Die getrockneten Sonnenblütenblätter werden fein im Mörser zerstoßen, bis sie ein feines Pulver ergeben, das vorsichtig eingerührt wird. Danach folgt das zerstoßene Blutkraut, welches dem Trank seine charakteristische, tiefrote Färbung verleiht. Der kritischste Moment ist die Beigabe des Taniumstaubs. Die magische Komponente muss unter ständigem Rühren in einer Acht-Bewegung hinzugefügt werden, um die heilenden Energien zu binden. Zum Schluss wird das Destillat durch ein feines Seidentuch gefiltert und in eine Phiole abgefüllt. 

Zubereitungszeit: Aktive Zeit: 3 Stunden (Konstantes Rühren und Temperaturkontrolle).

Ruhezeit: 12 Stunden in völliger Dunkelheit.

Menge: 3 Anwendungen (reicht für 3 separate Heilungen). 

Flüssigkeitsmenge: 15 Oth pro Anwendung (Gesamt 45 Oth). 

Gewicht: ca. 50 Oth. 

Wirkung: Heilung: Bei der Anwendung werden 1W3+3 Wunden geheilt.

    \paragraph{\trans{ingredients_label}}
    \begin{denseitemize}
        \item \textbf{Gemeines Blutkraut} (10 Oth), \textit{Zutaten}.
        Das Blutkraut ist ein bodennahes Gewächs, das vor allem durch seine fleischigen, tiefroten Blätter auffällt.

Die feinen Adern auf der Blattoberfläche leuchten in einem kräftigen Scharlachrot, fast so, als würde echtes Blut durch sie pulsieren. Wenn man ein Blatt zerreibt, tritt ein klebriger, süßlich riechender Saft aus, der die Finger tagelang verfärbt. Es ist keine magische Pflanze im klassischen Sinne. Sie bezieht ihre Kraft aus dem eisenreichen Boden und dem fahlen Licht der dichten Wälder.

Ihr findet das Blutkraut meist dort, wo es schattig und feucht ist. Es bevorzugt den Fuß alter Eichen oder die unmittelbare Nähe von verrottendem Unterholz in tiefen Wäldern. Ein unaufmerksamer Reisender hält es oft für gewöhnlichen Purpur-Ampfer, doch ein geschulter Alchemist achtet auf die kleinen, perlenartigen Tautropfen, die sich stets an den Rändern der Blätter sammeln.

Man erntet es am besten in den frühen Morgenstunden, bevor die Sonne das Blätterdach durchbricht. Man schneidet nur die äußeren Blätter ab, um die Wurzel zu schonen.
        \item \textbf{Taniumstaub} (1 Oth), \textit{Zutaten}.
        Tanium ist ein dunkles, kristallines Element von außergewöhnlicher Härte. Im Rohzustand findet man es oft als tiefschwarze Adern in uraltem Gestein, meist in Gebieten mit hoher natürlicher Magiekonzentration.

Es ist so hart, dass herkömmliche Mörser daran zerbrechen; nur Werkzeuge aus gehärtetem Diamant oder magisch verstärkte Mahlwerke können es zu feinem Staub zermahlen.

Tanium fungiert als perfekter Speicher für arkane Energie. Doch Vorsicht: Es hat kein Sättigungsgefühl! Wird es mit zu viel Magie gesättigt oder durch unsaubere Alchemie instabil, entlädt es sich in einer magischen Explosion.
        \item \textbf{Sonnenblüte} (5 Oth), \textit{Zutaten}.
        Die Sonnenblüte ist der Inbegriff von Beständigkeit. Mit ihrem kräftigen, rauen Stängel und dem stolzen, goldgelben Kranz aus Blütenblättern folgt sie unermüdlich dem Lauf des Sonnenwagens über das Firmament von Tiraka.

Ihr Kern ist gefüllt mit nahrhaften, ölhaltigen Samen, doch für Alchemisten sind es vor allem die leuchtenden Randblätter, die den Wert ausmachen. Sie speichern die reine, sanfte Wärme des Tages, ohne dabei die gefährliche Hitze des Feuers oder die Unberechenbarkeit der Magie in sich zu tragen.

Man findet sie fast überall in den fruchtbaren Ebenen Tirakas, besonders auf den sonnenverwöhnten Hügeln rund um Asgoran oder in den Gärten der Bauern im Hinterland. Sie liebt offene Flächen und tiefen, schwarzen Boden. Es ist keine seltene Pflanze, nein, aber eine, die Pflege braucht. Die wilden Varianten in den Heiden sind oft kleiner, aber ihre Essenz ist konzentrierter als die der gezüchteten Prachtexemplare.

Die Blätter sollten zur Mittagsstunde gezupft werden, wenn die Sonne am höchsten steht und die Blüte ihre volle Pracht entfaltet hat. Man trocknet sie flach liegend auf Leinentüchern an einem luftigen Ort, niemals im direkten Ofenfeuer!
        \item \textbf{Wasser} (45 Oth), \textit{Zutaten}.
        Kaltes klares Wasser.
    \end{denseitemize}


    \subsection*{Sud der seichten Ermächtigung}

    \begin{tabular}{@{}ll@{}}
        \textbf{\trans{difficulty_label}:} & Einfach \\
        \textbf{\trans{yield_label}:} & 2 \\
    \end{tabular}

    Zubereitung:
Der Schnaps wird mit der Frost-Flechte angesetzt, bis er eine bläuliche Färbung annimmt. Das Blut der Flugechse wird erst nach der Verfärbung untergerührt. Der Sud muss mindestens 5 Stunden an einem kühlen Ort ziehen (Keller, Höhle). Danach in die Phiole des Flugechsenblutes einfüllen, für den besten Effekt!

Zubereitungszeit: 2 Stunden

Menge: 2 Anwendungen, 2x200 Oth Flüssigkeit, 2x200 Oth Gewicht

Wirkung:
Der Spieler fügt seinem Würfelpool für die Dauer von 1W6 Minuten einen W6 hinzu. Dies gilt für alle Attributs-, Fertigkeits-, Kampf-, Magie-, Wissenswürfe, usw.
Nach Ablauf der Wirkungszeit ist der Spieler leicht reizbar und neigt zu Streitereien.

    \paragraph{\trans{ingredients_label}}
    \begin{denseitemize}
        \item \textbf{Phiole Flugechsenblut} (2 Phiole), \textit{Zutaten}.
        Meistens extrahiert aus den Adern der Flugechsen, die von den O'Gru domestiziert wurden. Es unterscheidet sich nicht vom Blut der in der Wildnis lebenden Exemplare.

Eine Phiole enthält 50 Oth Flugechselnblut.
        \item \textbf{Frostflechte} (100 Oth), \textit{Zutaten}.
        Ein zähes Gewächs, das ausschließlich in den nördlichen Steppen und den kältesten Regionen der Gebirge Tirakans gedeiht.

Sie wächst bevorzugt an der Baumgrenze und auf kargen, windgepeitschten Felsvorsprüngen, wo die Temperaturen selbst im Sommer kaum über den Gefrierpunkt steigen. Sie ist bei zwergischen Prospektoren aus dem Norden gut bekannt.

Die Flechte selbst besitzt keine eigene Magie, aber sie hat eine extreme kältebindende Eigenschaft. Sie absorbiert die arktische Kälte ihrer Umgebung und hält sie für bis zu 3 Tage gespeichert, wenn die Umgebungstemperatur nicht über 30° steigt. Innerhalb dieses Zeitraums beträgt die Temperatur der Pflanze weiterhin die des letzten Standorts.

In der Alchemie dient sie daher als Katalysator der Stabilisierung, der starke, flüchtige Substanzen (wie Blut oder hochprozentigen Alkohol) in einen Zustand des Kälteschocks versetzt.

Die Frostflechte erscheint als ein unauffälliges, dichtes Geflecht in Tiefblau oder Weißgrau. Sie liegt wie ein verkrusteter Teppich auf den Steinen und ist bei oberflächlicher Betrachtung kaum von eisüberzogenen Felsen oder von gefrorenem Moos zu unterscheiden. Sie besitzt keine Blätter und keine Blüten. 

Sie ist oft ein Handelsartikel in den südlichen Königreichen der Menschen, da sie dort nicht wächst, aber für einfache Heil- und Stärkungstränke unerlässlich ist.
        \item \textbf{Rübenschnaps} (400 Oth), \textit{Nahrung / Proviant}.
        Der Rübenschnaps ist ein typisches Erzeugnis der menschlichen Ackerländer in den gemäßigteren Zonen Tirakans, insbesondere in den westlichen und zentralen Königreichen der Menschen.

Er wird aus fermentierten und destillierten Zuckerrüben hergestellt.

Es handelt sich um einen billigen, aber hochprozentigen Fusel. Er dient in der Alchemie lediglich als Lösungsmittel und Hitzequelle, um die aggressiven oder flüchtigen Eigenschaften anderer Substanzen (wie die Galle der Flugechse) zu extrahieren. Aufgrund seiner Reinheit und seines Mangels an komplexen Inhaltsstoffen ist er ideal als einfache Basis für Massenträger.

Der Rübenschnaps ist meist klar oder leicht gelblich-trüb und riecht stechend nach Ethanol und einer unterschwelligen, erdigen Süße. Er brennt beim Trinken und hinterlässt einen starken, unangenehmen Nachgeschmack.

Der Rübenschnaps symbolisiert die Ausdauer und Pragmatik der Menschen in Tirakan. Während die Elfen ihr Lebendes Wasser und die Zwerge ihr Tiefen-Salz haben, verlassen sich die Menschen auf einfache, schnell verfügbare Lösungen.

Er ist ein Massenartikel und ein wichtiger Handelsgegenstand in den Grenzregionen und den Söldnerlagern. Ein Großteil des Preises einfacher Tränke entfällt auf die Destillation und den Transport, nicht auf die Zutat selbst.

Für die hochrangigen Alchemisten in den Akademien ist der Rübenschnaps ein Zeichen des Dilettantismus; sie bevorzugen raffiniertere, weniger aggressive Lösungsmittel. Die Dorfalchemisten hingegen nutzen ihn wegen seiner Effizienz und Verfügbarkeit.
    \end{denseitemize}


    \subsection*{Beturias ewige Ruhe}

    \begin{tabular}{@{}ll@{}}
        \textbf{\trans{difficulty_label}:} & Mittel \\
        \textbf{\trans{yield_label}:} & 1 \\
    \end{tabular}

    Zubereitung:
Der Brauprozess beginnt mit dem Kochen von Wasser, Baldrian und Monddornbeeren. Die Mischung muss zunächst auf die Hälfte ihrer ursprünglichen Flüssigkeitsmenge eingekocht werden. Nach dem Entfernen von der Hitze wird der Nachtschiefer-Staub langsam eingerührt, bis der Trank eine tiefschwarze Farbe annimmt. Der Sud muss anschließend einen ganzen Tag (24 Stunden) in einem versiegelten Gefäß reifen.

Zubereitungszeit: 1 Stunde, 24 Stunden bis zur Fertigstellung

Menge: 1 Anwendung, 300 Oth Flüssigkeit, 320 Oth Gewicht

Wirkung:
Nach der Einnahme fällt der Charakter sofort in einen 24-stündigen Scheintod-Zustand, wobei seine Körperfunktionen drastisch gesenkt werden und er 2W6 Wunden heilt. Um aus diesem beinahe komatösen Schlaf zu erwachen, benötigt der Charakter vier erfolgreiche Würfe auf Resistenz in Folge. Da Schütteln, laute Geräusche oder Angriffe ihn nicht wecken, kann er nur durch zwei erfolgreiche Würfe auf Untersuchen als lebend identifiziert werden; andernfalls wird er für tot gehalten.

    \paragraph{\trans{ingredients_label}}
    \begin{denseitemize}
        \item \textbf{Baldrian (Valeriana officinalis)} (4 Bund), \textit{Zutaten}.
        Baldrian hilft bei Einschlafstörungen und Unruhe. Hopfen und Melisse steigern die Wirkung des Baldrians und verbessern den Geschmack.
        \item \textbf{Nachtschiefer-Staub} (120 Oth), \textit{Zutaten}.
        Der Nachtschiefer ist kein gewöhnliches Mineral, sondern die fein zermahlene Essenz uralter, tiefschwarzer Gesteinsadern.

Dieses Material wird nicht geschürft, sondern muss aus dem Herz von Bergwerken gewonnen werden, die seit Äonen kein Sonnenlicht mehr gesehen haben, oft nur zugänglich durch enge Stollen in den kältesten, abgelegensten Gebirgszügen.

Man findet ihn ausschließlich in tiefen, stillgelegten Minen oder unterirdischen Krypten, wo das Gestein über Jahrtausende hinweg durch den konstanten Druck und die absolute Kälte des Erdinneren verdichtet wurde. Die Adern schimmern leicht, wenn sie im Schein einer Fackel erfasst werden. Ein Zeichen ihrer beinahe unnatürlichen Reinheit.
        \item \textbf{Wasser} (500 Oth), \textit{Zutaten}.
        Kaltes klares Wasser.
    \end{denseitemize}


    \subsection*{Elixier des saften Schlummers}

    \begin{tabular}{@{}ll@{}}
        \textbf{\trans{difficulty_label}:} & Einfach \\
        \textbf{\trans{yield_label}:} & 2 \\
    \end{tabular}

    Zubereitung:
Zuerst wird das Wasser in einem geeigneten Gefäß erhitzt und zum Kochen gebracht. Sobald es sprudelnd kocht, wird der Baldrian hinzugefügt und die Mischung muss für fünf Minuten kochen, damit sich die Wirkstoffe freisetzen können. Anschließend wird das Gefäß vom Feuer genommen und der Sud für weitere dreißig Minuten ziehen gelassen. Abschließend wird die fertige Flüssigkeit abgeseiht und in ein sauberes Gefäß gefüllt. Honig für den Geschmack hinzufügen.

Menge: 2 Anwendung, 2x250 Oth Flüssigkeit, 2x250 Oth Gewicht

Wirkung:
Nach dem Konsum fällt der Charakter innerhalb von zehn Minuten in einen festen Schlaf, der mindestens fünf Stunden anhält. Wird der Charakter während dieser Zeit gestört oder versucht man, ihn zu wecken, muss er einen Wurf auf Resistenz ablegen. Bei einem Erfolg erwacht der Charakter augenblicklich. Bei einem Misserfolg bleibt der Schlaf erhalten, und der Wurf kann bei der nächsten Störung oder einem erneuten Weckversuch wiederholt werden. Wichtig ist jedoch: Wird der Charakter geschüttelt, angegriffen oder sehr lauten Geräuschen ausgesetzt, erwacht er sofort.

    \paragraph{\trans{ingredients_label}}
    \begin{denseitemize}
        \item \textbf{Baldrian (Valeriana officinalis)} (1 Bund), \textit{Zutaten}.
        Baldrian hilft bei Einschlafstörungen und Unruhe. Hopfen und Melisse steigern die Wirkung des Baldrians und verbessern den Geschmack.
        \item \textbf{Wasser} (500 Oth), \textit{Zutaten}.
        Kaltes klares Wasser.
    \end{denseitemize}


    \subsection*{Elixier der elfischen Macht}

    \begin{tabular}{@{}ll@{}}
        \textbf{\trans{difficulty_label}:} & Mittel \\
        \textbf{\trans{yield_label}:} & 2 \\
    \end{tabular}

    Zubereitung:
Der Tran des Schillerwals wird mit dem Wasser zusammen in einer Glasschale langsam auf Körpertemperatur erwärmt. Er darf niemals sieden, da sonst der Schimmer (die magische Essenz) verfliegt. Das Steinsalz wird hineingerieselt, bis die Flüssigkeit milchig wird.
Die Silberorchideen-Blätter werden nun in das warme Öl eingelegt (Mazeration). Sie müssen dort 11 Stunden ruhen und langsam ihre Säfte an das Öl abgeben.
Am Ende wird das Elixier durch ein Seidentuch gefiltert. Die zurückbleibende Flüssigkeit sollte nun das charakteristische perlmuttartige Leuchten des Wals angenommen haben.

Zubereitungszeit: 1 Stunde aktive Zeit (für die genaue Temperaturkontrolle) und 11 Stunden Ruhezeit. 

Menge: 2 Anwendungen, 2x25 Oth Flüssigkeit, 2x25 Oth Gewicht

Wirkung:
Der Spieler fügt seinem Würfelpool für die Dauer von 2W6 Minuten  2W6 hinzu.
Nach Ablauf erleidet der Charakter den Zustand geschockt 2 und ist 2W6 Minuten Betäubt, da der Stoffwechsel nach dem Abklingen des hochpotenten Öls abrupt herunterfährt. Elfische Charaktere erleiden diese Nebenwirkung nicht.

    \paragraph{\trans{ingredients_label}}
    \begin{denseitemize}
        \item \textbf{Silberorchidee} (3 Blütenblätter), \textit{Zutaten}.
        Die Silberorchidee gilt als die unangefochtene und trügerische Königin der südlichen Flora. Sie ist ein botanisches Wunderwerk, das ebenso schön wie tödlich für jene ist, die sich von ihrem Glanz täuschen lassen.

Ihre Blätter sind nicht grün, sondern besitzen eine dunkelgraue, fast metallische Färbung, die im Mondlicht wie poliertes Silber glänzt. Die Adern pulsieren schwach in einem blassen Violett, wenn Magie in der Nähe ist. Die Blüte selbst ist groß und kelchförmig, mit schneeweißen Blütenblättern, die am Rand rasiermesserscharf sind. Doch das Verstörendste an ihr ist nicht ihre Schönheit, sondern ihre Mobilität: Die Pflanze reckt sich auf freiliegenden, muskulösen Wurzeln empor, mit denen sie in der Lage ist, langsam kriechend über den Boden zu wandern.

Die Silberorchidee ist fast ausschließlich im tiefen Süden Tirakans heimisch, jenseits der Eisernen Berge. Man findet sie in den weiten, grünen Steppen und an Flussläufen, oft im Schatten der dortigen Flora getarnt. Sie wächst oft in beunruhigender Nähe zu den riesenhaften Wesen aus Stein, die in den Pässen hausen.

Nähert man sich der Pflanze, stößt sie eine glitzernde Wolke aus feinem, silbernem Staub aus. Dies ist kein harmloser Pollenflug, sondern ein tödlicher Angriff. Wer den Staub einatmet, wird von schwerem Husten und Atemnot ergriffen. Binnen Augenblicken bilden sich schwarze Pocken auf der Haut, und das Opfer fällt in eine tiefe, todesähnliche Ohnmacht. Ist das Opfer wehrlos, sondert die Pflanze ein ätzendes Sekret aus ihren Wurzeln ab, um die Beute langsam zu zersetzen und aufzunehmen.

Trotz dieser Gefahren wird sie gejagt, denn sie ist ein mächtiger Potenzierer. In der Alchemie wird der extrahierte Nektar genutzt, um die Wirkung anderer Tränke ins Extreme zu steigern. Doch die Verarbeitung ist riskant: Ein Fehler in der Destillation führt dazu, dass die magische Energie den Körper überlädt (analog zum Silbertod), was oft tödlich endet.

In den alten Legenden heißt es, die Silberorchideen seien entstanden, als im Ersten Zeitalter das Blut eines gefallenen Sternengottes auf die Erde tropfte. Die Elfen hingegen nennen die Blume Verräterschmuck und glauben, dass sie dort wächst, wo die Realität Risse bekommen hat und Chaos in die Welt sickert.

Sieht aus wie Schmuck, kriecht wie eine Spinne und ist mehr wert als mein Haus. Aber pass auf, Junge: Wenn du das Glitzern siehst, halte den Atem an und renn. Bevor du merkst, dass du krank wirst, bist du schon ihr Dünger. – Randnotiz in den Aufzeichnungen des Kräutersammlers 'Drei-Finger-Hannes'
        \item \textbf{Tran des Schillerwals} (1 Phiole), \textit{Zutaten}.
        Der Tran, das Fettgewebe unter der Haut, ist der Grund, warum diese majestätischen Tiere gejagt werden. Es ist kein gewöhnliches Fett, sondern ein  Speicher magischer Energie.

Der Tran ist im rohen Zustand eine zähe, geleeartige Masse, die schwach bläulich leuchtet. Nach der Veredelung (Schmelzen und Filtern) wird er zu einem klaren, öligen Elixier, das wie flüssiges Perlmutt schlieren zieht.
Im Gegensatz zum ranzigen Tran gewöhnlicher Wale riecht der Tran des Schillerwals frisch, salzig und leicht metallisch (wie die Luft vor einem Gewitter).

Er ist das beste bekannte Mittel, um flüchtige Magie in Tränken zu binden (siehe Elixier der Elfenwacht).

In Lampen verbrannt, gibt er ein Licht ab, das niemals rußt und unsichtbares sichtbar machen kann. Waffenöle aus diesem Tran lassen Geister verletzen.
    \end{denseitemize}


    \subsection*{Lachsunder Seemagen}

    \begin{tabular}{@{}ll@{}}
        \textbf{\trans{difficulty_label}:} & Einfach \\
        \textbf{\trans{yield_label}:} & 3 \\
    \end{tabular}

    Zubereitung:
Nehmt Euren Kleinen Kessel und füllt ihn mit klarem Quellwasser. Zerstoßt die getrocknete Beifußpflanze und die Algen mit Mörser und Stößel, bis ein grober Brei entsteht. Nun gebt Ihr frisch geriebenen Ingwer hinzu. Lasst das Ganze kurz aufkochen und sodann bei mäßiger Hitze sieden, bis sich die Flüssigkeit dunkelgrün färbt. Seiht den Sud durch ein Leinentuch ab und füllt ihn in eine Phiole.

Zubereitungszeit:
20 Minuten aktive Brauzeit, keine Ruhezeit erforderlich.

Menge:
3 Anwendung (Flüssigkeit: 30 Oth, Gesamtgewicht 90 Oth).

Wirkung:
Heilt sofort alle Mali die durch Seekrankheit entstanden sind (kann auch präventiv bis zu 2 Stunden vor der Seereise eingenommen werden) und nach 2 Stunden eine Wunde.

    \paragraph{\trans{ingredients_label}}
    \begin{denseitemize}
        \item \textbf{Beifuß (Artemisia vulgaris)} (1 Bund), \textit{Zutaten}.
        Eine Beifuß Pflanze. Die Spitzen des Sprosses werden verwendet um den Gallenfluss und die Verdauung in Gang zu bringen.
        \item \textbf{Lachsunder Seetang (Alga asgorania)} (10 Oth), \textit{Zutaten}.
        Eine zähe, dunkelgrüne bis bräunliche Meeresalge mit breiten, ledrigen Blättern, die an den Küsten von Lachsund im kalten Meerwasser gedeiht.

Sie schmeckt extrem salzig und verströmt einen intensiven Duft nach frischer Gischt und Jod.

In der alchemistischen Naturkunde dient er als Stabilisator für den Magen. Seine Inhaltsstoffe beruhigen den Gleichgewichtssinn und binden Giftstoffe im Verdauungstrakt. Für den alltäglichen Verzehr ist er nicht geeignet

Das Gewächs wird in großen Mengen an den felsigen Küsten von Asgoran angespült oder gezielt von den maritimen Sammlern der Siederei Falke aus dem Meer gefischt, getrocknet und verarbeitet. Auf den Fischmärkten von Lachsund ist er körbeweise erhältlich.
        \item \textbf{Zitronen-Ingwer (Zingiber meridionale)} (5 Oth), \textit{Zutaten}.
        Eine knollige, faserige Wurzel von hellbrauner Farbe, deren saftiges, gelbes Fleisch beim Aufschneiden eine brennende Schärfe und ein intensiv frisches Zitrusaroma freisetzt.

In der Alchemie gilt er als feuriger, die Durchblutung anregender Vitalisator. Er wird genutzt, um die Säfte des Körpers in Wallung zu bringen, Magenbeschwerden zu lindern und die Aufnahme von Heilelixieren im Blut drastisch zu beschleunigen.

Die Pflanze liebt feuchte Wärme und gedeiht prächtig in den fruchtbaren, flussreichen Ländern der O'grut oder wird über die weitreichenden Handelsnetzwerke der Menschen aus den südlichen Oasen nach Meridian importiert.
        \item \textbf{Wasser} (90 Oth), \textit{Zutaten}.
        Kaltes klares Wasser.
    \end{denseitemize}


    \subsection*{Panthomeischer Wachmacher}

    \begin{tabular}{@{}ll@{}}
        \textbf{\trans{difficulty_label}:} & Mittel \\
        \textbf{\trans{yield_label}:} & 2 \\
    \end{tabular}

    Zubereitung:
Zunächst müsst ihr die harten Panthomeischen Nadelzapfen in einem Mörser zu feinem Staub zermahlen. Dieser Staub wird in klarem Wasser über einer blauen, extrem heißen Flamme aufgekocht, bis das Wasser sich smaragdgrün färbt. Nun müsst ihr frisches Blutkraut hinzugeben, um die bitteren Stoffe der Zapfen zu neutralisieren. Die Kamille, leicht zerstoßen, verleiht dem Trunk seine energetisierende Wirkung.

Zubereitungszeit:
2 Stunden hochkonzentriertes Kochen, gefolgt von einer halben Nacht Ziehzeit im kühlen Keller.

Menge:
2 Anwendungen. 200 Oth Flüssigkeit, Gesamtgewicht 200 Oth.

Wirkung:
Beseitigt temporär Schlafmangel. Der Konsument ignoriert alle Abzüge durch Müdigkeit und erhält zudem einen satten Bonus von +2 Würfeln auf Wahrnehmung und Wachsamkeit für 1W6 Stunden. Sobald die Wirkung jedoch abklingt, stürzt der Körper unweigerlich in einen tiefen, traumlosen Schlaf.

    \paragraph{\trans{ingredients_label}}
    \begin{denseitemize}
        \item \textbf{Panthomeischer Nadelzapfen} (2 Stück), \textit{Zutaten}.
        Panthomeische Nadelzapfen sind die verholzten, harzreichen Fruchtstände der seltenen Panthomeischen Gebirgskiefer.

Diese Zapfen sind von einer harten, fast metallisch wirkenden Schuppenschicht überzogen, die im Inneren einen tiefvioletten, hochgradig klebrigen Balsam einschließt. Sie besitzen keine inhärente magische Energie, sind jedoch in der Alchemie für ihre stark konservierenden, strukturfestigenden und wärmespenden Eigenschaften bekannt.

Diese Zapfen wachsen ausschließlich in den eisigen, sturmdurchpeitschten Höhen der Panthomeischen Gipfel, weit oberhalb der Baumgrenze im unwegsamen Hochland.

Die Ernte ist beschwerlich, da die Zapfen nur von den Jägern oder den Hil-JabalJarh' geschlagen werden müssen.

Gelegentlich bringen zwergische Gebirgskarawanen vereinzelte Exemplare auf die Märkte in ganz Tirakan.
        \item \textbf{Gemeines Blutkraut} (20 Oth), \textit{Zutaten}.
        Das Blutkraut ist ein bodennahes Gewächs, das vor allem durch seine fleischigen, tiefroten Blätter auffällt.

Die feinen Adern auf der Blattoberfläche leuchten in einem kräftigen Scharlachrot, fast so, als würde echtes Blut durch sie pulsieren. Wenn man ein Blatt zerreibt, tritt ein klebriger, süßlich riechender Saft aus, der die Finger tagelang verfärbt. Es ist keine magische Pflanze im klassischen Sinne. Sie bezieht ihre Kraft aus dem eisenreichen Boden und dem fahlen Licht der dichten Wälder.

Ihr findet das Blutkraut meist dort, wo es schattig und feucht ist. Es bevorzugt den Fuß alter Eichen oder die unmittelbare Nähe von verrottendem Unterholz in tiefen Wäldern. Ein unaufmerksamer Reisender hält es oft für gewöhnlichen Purpur-Ampfer, doch ein geschulter Alchemist achtet auf die kleinen, perlenartigen Tautropfen, die sich stets an den Rändern der Blätter sammeln.

Man erntet es am besten in den frühen Morgenstunden, bevor die Sonne das Blätterdach durchbricht. Man schneidet nur die äußeren Blätter ab, um die Wurzel zu schonen.
        \item \textbf{Kamille (Matricaria recutita)} (1 Bund), \textit{Zutaten}.
        Kamille zählt zu den ältesten Heilpflanzen und wurde bereits im Mittelalter angewandt. Die Blüten haben eine heilende und beruhigende Wirkung. Äußerlich kann Kamille bei Entzündungen des Zahnfleisches, der Haut oder der Schleimhaut angewendet werden. Innerlich eingenommen wirkt sie bei Magen-Darm-Erkrankungen. Spülen und die Inhalation sind ebenfalls weit verbreitet.
        \item \textbf{Wasser} (130 Oth), \textit{Zutaten}.
        Kaltes klares Wasser.
    \end{denseitemize}


    \subsection*{Schlangenöl}

    \begin{tabular}{@{}ll@{}}
        \textbf{\trans{difficulty_label}:} & Einfach \\
        \textbf{\trans{yield_label}:} & 1 \\
    \end{tabular}

    Zubereitung:
Zuerst wird das Fett (irgendeines Hauptsache fettig!) in einem kleinen Eisentiegel bei mäßiger Hitze ausgelassen. Achtet darauf, dass es nicht verbrennt. Sobald das Fett flüssig ist, rührt Ihr den Absud aus Minze unter. Diese dient allein dazu, den stechenden Geruch des Fetts zu überdecken und der Tinktur eine medizinisch-seriöse grüne Färbung zu verleihen. Ein Schuss Rübenschnaps sorgt dafür, dass die Mischung im Hals brennt. Den nur was brennt, das hilft auch. Zum Schluss wird die Emulsion kräftig geschüttelt, bis sie trüb und geheimnisvoll aussieht.

Zubereitungszeit: 30 Minuten aktive Rührarbeit, 1 Stunde zum Abkühlen und Absetzen.

Menge: 1 Anwendung (Flüssigkeit: 5 Oth / Gewicht: 6 Oth inklusive Fläschchen).

Wirkung: Stellt bei Anwendung sofort 1 Wunde wieder her. Der Patient fühlt sich kurzzeitig belebt, leidet aber unter dem schrecklichen Nachgeschmack.

    \paragraph{\trans{ingredients_label}}
    \begin{denseitemize}
        \item \textbf{Fett} (10 Oth), \textit{Zutaten}.
        Es ist Fett. Pflanzlichem oder tierischem Ursprung. Es dient zum braten, zum Verfeinern von Speisen oder zum Schmieren.
        \item \textbf{Minze} (1 Bund), \textit{Zutaten}.
        Die Minze ist in Tirakan weit verbreitet, doch unter jenen, die sich mit der Kräuterkunde befassen, gilt sie als unverzichtbares Grundnahrungsmittel für Geist und Körper.

Vorkommen \& Wuchs:
Man findet die Minze meist an feuchten, halbschattigen Orten. Sie wuchert an den Ufern von Bachläufen, in verwunschenen Waldlichtungen oder in den Kräutergärten weiser Heilerinnen. 

Sie gilt als magenberuhigend, wirkt kühlend und lindert Halsschmerzen. Sie findet verwendung in Salben, Tinkturen und Tees. Wird auch oft genutzt, um den strengen Geschmack von Wildfleisch zu mildern oder um billigem Dünnbier eine frische Note zu verleihen.
        \item \textbf{Wasser} (10 Oth), \textit{Zutaten}.
        Kaltes klares Wasser.
    \end{denseitemize}


    \subsection*{Sud des unkontrollierten Fleißes (Goblin Variante)}

    \begin{tabular}{@{}ll@{}}
        \textbf{\trans{difficulty_label}:} & Einfach \\
        \textbf{\trans{yield_label}:} & 1 \\
    \end{tabular}

    Zubereitung:
Die getrockneten Blätter des Gildenkrauts werden grob mit den Fingern zerfetzt und zusammen mit billigem Rübenschnaps in einen unbgereinigten Eisentiegel geworfen. Unter lautem, rhythmischem Klatschen wird das Gemisch einmal kurz und heftig zum Sieden gebracht, bis ein stechender, nussiger Dampf aufsteigt. Sofort vom Feuer nehmen und warm in eine lederne Blase abfüllen.

Zubereitungszeit:
15 Minuten aktive Arbeit (keine Ruhezeit).

Menge: 1 Anwendung (10 Oth Flüssigkeit, Gewicht 12 Oth).

Wirkung: Dieser Sud wirkt nur auf Individuen der Goblinvölker. Trinkt ein anderes Wesen diesen Trank, fühlt dieses sich etwas munterer und ist voller Tatendrang.

Der Goblin jedoch verfällt für 1W6x10 Minuten in einen Arbeitswahn. Er erhält +2 Würfel auf alle handwerklichen oder körperlichen Proben (z.B. Mechanik, Bergbau, Klettern). Allerdings sind alle feingliedrigen oder logischen Proben (Logik, Täuschen und Betrügen) in dieser Zeit um +3 erschwert, da der Geist völlig unkoordiniert hin und her springt. Er kann nicht an sich halten, kann kaum aufgehalten werden, außer von starken Armen oder Fesseln.

    \paragraph{\trans{ingredients_label}}
    \begin{denseitemize}
        \item \textbf{Gildenkraut (Vigilia mercatoris)} (5 Oth), \textit{Zutaten}.
        Das Gildenkraut ist eine zähe, beinahe unscheinbare Pflanze, die nur dem geschulten Auge eines Kundigen auffällt.

Sie zeichnet sich durch gezackte, ledrige Blätter von tiefgrüner Färbung aus, deren dicke Blattadern einen leichten rötlichen Schimmer aufweisen. Ihre Blüten sind winzig, blassgelb und schließen sich sofort, sobald die Sonne ihren Zenit überschreitet.

Das Kraut ist vollkommen mundan und besitzt nicht den geringsten Funken Arkana. Sein Wert liegt vielmehr in der extrem hohen, natürlichen Konzentration an Alkaloiden, vor allem reinem Koffein. Zerreibt man die harten Blätter zwischen den Fingern, verströmen sie einen scharfen, erdigen Duft, der an stark geröstete Nüsse erinnert.

Lage und Lebensraum:
Dieses Gewächs ist verhältnismäßig selten, da es ganz bestimmte, widrige Umstände benötigt, um seine anregenden Säfte zu bilden. Es gedeiht ausschließlich in kargen, mineralreichen Böden und meidet das feuchte Moos tiefer Wälder.

Man findet das Gildenkraut vornehmlich in den felsigen Spalten uralter Handelsrouten, im bröckelnden Steinmörtel verlassener Gilden-Außenposten oder an den steilen, sonnengegerbten Hängen des Toraner Hochlands. Es braucht den harten Fels und viel Wind, um nicht zu verkümmern.
        \item \textbf{Rübenschnaps} (10 Oth), \textit{Nahrung / Proviant}.
        Der Rübenschnaps ist ein typisches Erzeugnis der menschlichen Ackerländer in den gemäßigteren Zonen Tirakans, insbesondere in den westlichen und zentralen Königreichen der Menschen.

Er wird aus fermentierten und destillierten Zuckerrüben hergestellt.

Es handelt sich um einen billigen, aber hochprozentigen Fusel. Er dient in der Alchemie lediglich als Lösungsmittel und Hitzequelle, um die aggressiven oder flüchtigen Eigenschaften anderer Substanzen (wie die Galle der Flugechse) zu extrahieren. Aufgrund seiner Reinheit und seines Mangels an komplexen Inhaltsstoffen ist er ideal als einfache Basis für Massenträger.

Der Rübenschnaps ist meist klar oder leicht gelblich-trüb und riecht stechend nach Ethanol und einer unterschwelligen, erdigen Süße. Er brennt beim Trinken und hinterlässt einen starken, unangenehmen Nachgeschmack.

Der Rübenschnaps symbolisiert die Ausdauer und Pragmatik der Menschen in Tirakan. Während die Elfen ihr Lebendes Wasser und die Zwerge ihr Tiefen-Salz haben, verlassen sich die Menschen auf einfache, schnell verfügbare Lösungen.

Er ist ein Massenartikel und ein wichtiger Handelsgegenstand in den Grenzregionen und den Söldnerlagern. Ein Großteil des Preises einfacher Tränke entfällt auf die Destillation und den Transport, nicht auf die Zutat selbst.

Für die hochrangigen Alchemisten in den Akademien ist der Rübenschnaps ein Zeichen des Dilettantismus; sie bevorzugen raffiniertere, weniger aggressive Lösungsmittel. Die Dorfalchemisten hingegen nutzen ihn wegen seiner Effizienz und Verfügbarkeit.
    \end{denseitemize}


    \subsection*{Transalbe der Falk's}

    \begin{tabular}{@{}ll@{}}
        \textbf{\trans{difficulty_label}:} & Mittel \\
        \textbf{\trans{yield_label}:} & 2 \\
    \end{tabular}

    Zubereitung:
Nehmt einen kleinen Kessel aus Eisen und erhitzt den Schillerwal-Tran bei mäßigem Feuer, bis er vollkommen flüssig und klar wird.

Gebt unter ständigem, rhythmischem Rühren den getrockneten und feinst zerriebenen Spitzwegerich hinzu, um der Salbe ihre wundschließende und hautstraffende Eigenschaft zu verleihen. Sobald die Mischung leicht bläulich schimmert, zieht den Kessel von den Kohlen und rührt das gereinigte Fett unter, während die Masse langsam in die Steinguttiegel abgefüllt wird und dort zu einer zähen Paste erstarrt.

Zubereitungszeit:
1 Stunde aktives Sieden und Rühren, 3 Stunden Ruhezeit zum Erstarren.

Menge:
2 Anwendungen, 70 Oth reine Salbenmasse.

Wirkung:
Das Auftragen erfordert 1 Minute Zeit. Gewährt dem Anwender für die nächsten 2 Stunden 1 Punkt Wundheilung (absorbiert eine beliebige Wunde) sowie die Schutzeinheiten N (Normaler Schutz) und B (Blutungsschutz) für den nächsten Kampf. Zudem ist der Anwender in dieser Zeit immun gegen die Effekte natürlicher Kälte.

    \paragraph{\trans{ingredients_label}}
    \begin{denseitemize}
        \item \textbf{Tran des Schillerwals} (50 Oth), \textit{Zutaten}.
        Der Tran, das Fettgewebe unter der Haut, ist der Grund, warum diese majestätischen Tiere gejagt werden. Es ist kein gewöhnliches Fett, sondern ein  Speicher magischer Energie.

Der Tran ist im rohen Zustand eine zähe, geleeartige Masse, die schwach bläulich leuchtet. Nach der Veredelung (Schmelzen und Filtern) wird er zu einem klaren, öligen Elixier, das wie flüssiges Perlmutt schlieren zieht.
Im Gegensatz zum ranzigen Tran gewöhnlicher Wale riecht der Tran des Schillerwals frisch, salzig und leicht metallisch (wie die Luft vor einem Gewitter).

Er ist das beste bekannte Mittel, um flüchtige Magie in Tränken zu binden (siehe Elixier der Elfenwacht).

In Lampen verbrannt, gibt er ein Licht ab, das niemals rußt und unsichtbares sichtbar machen kann. Waffenöle aus diesem Tran lassen Geister verletzen.
        \item \textbf{Spitzwegerich (Plantago lanceolata)} (1 Bund), \textit{Zutaten}.
        Die spitzen, schmalen Blätter des Spitzwegerich werden als Sirup oder auch als Tee bei Erkältungskrankheiten verwendet. Spitzwegerich kann auch zerkleinert und zerrieben auf Wunden oder Insektenstiche aufgetragen werden und entfaltet hier eine kühlende Wirkung. Die Pflanze wird auch bei Durchfall eingesetzt.
        \item \textbf{Fett} (20 Oth), \textit{Zutaten}.
        Es ist Fett. Pflanzlichem oder tierischem Ursprung. Es dient zum braten, zum Verfeinern von Speisen oder zum Schmieren.
    \end{denseitemize}


    \subsection*{Trank des Schutzes}

    \begin{tabular}{@{}ll@{}}
        \textbf{\trans{difficulty_label}:} & Einfach \\
        \textbf{\trans{yield_label}:} & 1 \\
    \end{tabular}

    Zubereitung:
Zuerst müsst Ihr das Quellwasser in einem Kupferkessel zum Sieden bringen, bis der Dampf die Fenster beschlägt. Gebt die getrockneten Beifußblätter hinzu. Sobald das Wasser eine tiefgrüne Färbung annimmt, rührt Ihr das Kiefernharz ein. Hier ist Vorsicht geboten: Rührt langsam und stets im Uhrzeigersinn, sonst klumpt das Harz und Ihr habt am Ende nur klebrigen Schleim statt flüssigen Schutzes! Zum Schluss wird die Alant-Wurzel fein zerstoßen und hinzugefügt, um die Vitalität zu binden. Filtert das Ganze durch ein feines Leinentuch in eine Phiole. 

Zubereitungszeit: 45 Minuten aktive Brauzeit, danach 1 Stunden Ruhezeit zum Abkühlen und Binden der Essenzen.

Menge:
1 Anwendung (ca. 8 Oth Flüssigkeit, Gewicht der Phiole inklusive Inhalt: 10 Oth).

Wirkung:
Bei der Anwendung erhält der Charakter 1W3 Boost auf seine Rüstung. Jeder dieser Boosts kann als Normaler Schutz (N) eingesetzt werden. Das bedeutet, ein Boost kann verbraucht werden, um einen normalen Treffer komplett zu verhindern.

    \paragraph{\trans{ingredients_label}}
    \begin{denseitemize}
        \item \textbf{Wasser} (50 Oth), \textit{Zutaten}.
        Kaltes klares Wasser.
        \item \textbf{Beifuß (Artemisia vulgaris)} (1 Bund), \textit{Zutaten}.
        Eine Beifuß Pflanze. Die Spitzen des Sprosses werden verwendet um den Gallenfluss und die Verdauung in Gang zu bringen.
        \item \textbf{Alant (Inula helenium)} (1 Bund), \textit{Zutaten}.
        Diese Heilpflanze aus dem Mittelalter ist in modernen Zeiten nicht mehr weit verbreitet. Ihre Anwendung verbessert die Verdauung, und ihr wird eine vorbeugende Wirkung gegen Darmkrebs zugeschrieben.
        \item \textbf{Kiefernharz} (10 Oth), \textit{Zutaten}.
        Das schützende Blut des Baumes. Wenn die raue Rinde einer Kiefer,sei es durch den Biss eines Wildtieres oder das ungeschickte Beil eines Holzfällers, verletzt wird, tritt diese zähe, goldgelbe Flüssigkeit aus.

Sie fließt langsam, beinahe träge, und erfüllt die Umgebung mit diesem unverwechselbaren, herben Duft.

Es wird genutzt, um spröde Korken zu versiegeln, Fackeln eine kräftige Flamme zu verleihen oder um einfache Wundsalben dickflüssig zu machen.

Sobald es an der Luft trocknet, wird es steinhart und schließt die Wunde des Baumes wie ein natürlicher Pfropfen. 

Wenn Ihr das Zeug erst einmal an den Händen habt, hilft nur noch ordentlich Fett oder sehr viel Geduld, um den Kleber wieder loszuwerden.
    \end{denseitemize}


    \subsection*{Elexir des Chronos}

    \begin{tabular}{@{}ll@{}}
        \textbf{\trans{difficulty_label}:} & Schwer \\
        \textbf{\trans{yield_label}:} & 1 \\
    \end{tabular}

    Zubereitung:
Bringt das reine Quellwasser in einem schweren Eisenkessel zum Sieden. Fügt den Obsidianstaub hinzu und rührt ununterbrochen, bis das Wasser eine trübe, graue Färbung annimmt. Nun gebt das getrocknete Gildenkraut und den Baldrian hinzu. Der Staub bindet die Pflanzenfasern und sinkt während des dreistündigen, schweren Siedens langsam zu Boden, während er die konzentrierten Alkaloide in der Flüssigkeit freisetzt.

Das Filtern muss hier durch ein dreifaches Leinentuch erfolgen. Verbleibt Obsidianstaub im Trank, liegt er dem Konsumenten schwer in den Eingeweiden, was 1W6 Stunden Erschöpfung verursacht. Alle Proben sind während dieser Zeit um +1 erschwert.

Zubereitungszeit:
3 Stunden aktive, hochkonzentrierte Arbeit. 12 Stunden Ruhezeit in absoluter Dunkelheit und Stille.

Menge: 1 Anwendung (ca. 15 Oth reine Essenz, Gesamtgewicht 15 Oth).

Wirkung:
Der Konsument tritt für die nächsten 6 Stunden in einen Zustand absoluter, übernatürlicher Hyper-Fokussierung.
Der Charakter erhält einen permanenten Bonus von +2 Würfeln auf alle Proben, die Logik, Gewissenhaftigkeit, Auffassungsgabe oder Wissen erfordern.

Zudem wird der Charakter für die Dauer der Wirkung vollkommen immun gegen den Zustand Benommenheit sowie gegen alle mentalen Beeinflussungsversuche oder Schlafzauber (da das Gehirn schlicht zu schnell arbeitet).

Sobald die 6 Stunden ablaufen, bricht das Nervensystem schlagartig zusammen. Der Charakter fällt augenblicklich für 1W6+2 Stunden in einen komatösen Tiefschlaf, aus dem er absolut nicht geweckt werden kann (selbst Schaden weckt ihn nicht). Er erwacht danach etwas verwirrt und kann sich nicht erinnern, was innerhalb der 6 stündigen Wirkungszeit passiert ist. 

Der Tiefschlaf bewirkt jedoch, dass der Charakter einen Rast-Wurf werfen darf.

    \paragraph{\trans{ingredients_label}}
    \begin{denseitemize}
        \item \textbf{Gildenkraut (Vigilia mercatoris)} (30 Oth), \textit{Zutaten}.
        Das Gildenkraut ist eine zähe, beinahe unscheinbare Pflanze, die nur dem geschulten Auge eines Kundigen auffällt.

Sie zeichnet sich durch gezackte, ledrige Blätter von tiefgrüner Färbung aus, deren dicke Blattadern einen leichten rötlichen Schimmer aufweisen. Ihre Blüten sind winzig, blassgelb und schließen sich sofort, sobald die Sonne ihren Zenit überschreitet.

Das Kraut ist vollkommen mundan und besitzt nicht den geringsten Funken Arkana. Sein Wert liegt vielmehr in der extrem hohen, natürlichen Konzentration an Alkaloiden, vor allem reinem Koffein. Zerreibt man die harten Blätter zwischen den Fingern, verströmen sie einen scharfen, erdigen Duft, der an stark geröstete Nüsse erinnert.

Lage und Lebensraum:
Dieses Gewächs ist verhältnismäßig selten, da es ganz bestimmte, widrige Umstände benötigt, um seine anregenden Säfte zu bilden. Es gedeiht ausschließlich in kargen, mineralreichen Böden und meidet das feuchte Moos tiefer Wälder.

Man findet das Gildenkraut vornehmlich in den felsigen Spalten uralter Handelsrouten, im bröckelnden Steinmörtel verlassener Gilden-Außenposten oder an den steilen, sonnengegerbten Hängen des Toraner Hochlands. Es braucht den harten Fels und viel Wind, um nicht zu verkümmern.
        \item \textbf{Baldrian (Valeriana officinalis)} (1 Bund), \textit{Zutaten}.
        Baldrian hilft bei Einschlafstörungen und Unruhe. Hopfen und Melisse steigern die Wirkung des Baldrians und verbessern den Geschmack.
        \item \textbf{Obsidian (Staub)} (5 Oth), \textit{Zutaten}.
        Obsidianstaub ist ein mundaner, mineralischer alchemistischer Grundstoff, der aus vulkanischem Glas gewonnen wird.

Er besitzt keinerlei innewohnende magische Kraft, ist jedoch aufgrund seiner extremen Härte, seiner scharfen Kanten im mikroskopischen Bereich und seiner enormen Hitzebeständigkeit ein unschätzbarer mechanischer Katalysator.

In der Alchemie wird er genutzt, um Unreinheiten und Bitterstoffe aus kochenden Suden zu binden und feine Pflanzenessenzen zu zwingen, sich abzusetzten.

Obsidian entsteht dort, wo feurige Magma rasch an der frischen Luft oder durch Wasser abkühlt. Man findet das vulkanische Glas in den tiefen, erstarrten Lavaseen der Unterwelt oder an den rauen, vulkanischen Bruchstellen des Hochlands.
Er kann in schweren Mörsern zu einem mehlfeinen, grauen Staub zerstoßen werden.

Die unangefochtenen Hauptlieferanten für den besten Obsidianstaub in Tirakan sind die Zwerge der Xordai. Insbesondere die Kolonie Kor-Vamak, die direkt in einem erstarrten Lavasee errichtet wurde, hat sich auf den Abbau und die Verarbeitung dieses Materials spezialisiert.

Das dort ansässige Syndikat exportiert den Staub über ihre gut bewachten Handelskarawanen bis an die Oberfläche.
    \end{denseitemize}


    \subsection*{Felsenfeuers Destillat}

    \begin{tabular}{@{}ll@{}}
        \textbf{\trans{difficulty_label}:} & Einfach \\
        \textbf{\trans{yield_label}:} & 1 \\
    \end{tabular}

    Zubereitung:
Die Zubereitung verlangt meridianische
Präzision. Zunächst wird die Eibischwurzel kalt angesetzt
und exakt zwei Stunden ziehen gelassen, bevor die Flüssig-
keit gesiebt und erhitzt wird. In er Zwischenzeit werden Beinwell und Kamille in einem Mörser zu  einem feinen Pulver zerrieben. Die erhitzte Eibischflüssigkeit wird mit dem Pulver vermengt und mehrfach destilliert, bis eine klare Esenz entstet.

Zubereitungszeit:
2 Stunden Ruhezeit für den Kaltaus-
zug der Eibischwurzel, gefolgt von 1 Stunde aktiver Zeit
an den Destillationsapparaturen.

Menge:
1 Anwendung. Ca. 200 Oth Flüssigkeit, Gesamt-
gewicht 800 Oth (inklusive Keramikflasche).

Wirkung:
Die konzentrierten Wirkstoffe beschleunigen die Regeneration der Zellen und heilen 1W3 Wunden.

    \paragraph{\trans{ingredients_label}}
    \begin{denseitemize}
        \item \textbf{Wasser} (200 Oth), \textit{Zutaten}.
        Kaltes klares Wasser.
        \item \textbf{Kamille (Matricaria recutita)} (50), \textit{Zutaten}.
        Kamille zählt zu den ältesten Heilpflanzen und wurde bereits im Mittelalter angewandt. Die Blüten haben eine heilende und beruhigende Wirkung. Äußerlich kann Kamille bei Entzündungen des Zahnfleisches, der Haut oder der Schleimhaut angewendet werden. Innerlich eingenommen wirkt sie bei Magen-Darm-Erkrankungen. Spülen und die Inhalation sind ebenfalls weit verbreitet.
        \item \textbf{Beinwell (Symphytum officinale)} (50), \textit{Zutaten}.
        Beinwell regt die Durchblutung an, Prellungen, Hämatome und Verstauchungen verschwinden schneller. Beinwell beschleunigt die Regeneration der Zellen.
        \item \textbf{Eibisch (Althaea officinalis)} (100), \textit{Zutaten}.
        Von dieser Heilpflanze wird die Wurzel verwendet. Diese wird kalt angesetzt und muss ungefähr zwei Stunden ziehen. Erst nach dem Ziehen wird die Flüssigkeit gesiebt und anschließend erhitzt. Die Stoffe bieten Schutz für die Schleimhäute und wirken reizlindernd. Eine hilfreiche Heilpflanze bei Magen-Darm-Problemen und Husten.
    \end{denseitemize}


    \subsection*{Nebel von Thenon}

    \begin{tabular}{@{}ll@{}}
        \textbf{\trans{difficulty_label}:} & Mittel \\
        \textbf{\trans{yield_label}:} & 1 \\
    \end{tabular}

    Zubereitung:
Nehmt eine saubere Glasphiole und beträufelt den Boden mit dem Sekret des Kinstarchel-Knochenextrakts, um die explosive Freisetzung bei Aufprall zu sichern.

Kocht parallel die getrockneten Spitzwegerichblätter in reinem Quellwasser, bis die Flüssigkeit eine sämige, schleimige Konsistenz annimmt. Fügt unter ständigem Rühren gegen den Uhrzeigersinn den gemahlenen Bernstein hinzu. Gebt das Harz hinzu, es bindet die Dämpfe. Versiegelt die Phiole sofort mit flüssigem Wachs, solange der Dampf noch heiß ist.

Zubereitungszeit:
3 Stunden aktive Alchemie, danach 1 Stunde Abkühlung im Dunkeln.

Menge:
1 Anwendung (Phiole fasst ca. 50 Oth Flüssigkeit)

Wirkung:
Erzeugt beim Zerbrechen augenblicklich eine stationäre Nebelwolke in einem Radius von 3 Metern, die für 5 Runden bestehen bleibt. Alle Kreaturen darin erhalten den Zustand Blind und jegliche Sicht ist stark eingeschränkt (Angriffe und Wahrnehmungsproben von außen nach innen und umgekehrt gelten als schwere Proben mit +2 MW).

    \paragraph{\trans{ingredients_label}}
    \begin{denseitemize}
        \item \textbf{Kinstarchel Sekret} (10 Oth), \textit{Zutaten}.
        Eine hoch reaktive Flüssigkeit, die aus den Knochen der possierlichen Kinstarchel hergestellt wird.
        \item \textbf{Spitzwegerich (Plantago lanceolata)} (1 Bund), \textit{Zutaten}.
        Die spitzen, schmalen Blätter des Spitzwegerich werden als Sirup oder auch als Tee bei Erkältungskrankheiten verwendet. Spitzwegerich kann auch zerkleinert und zerrieben auf Wunden oder Insektenstiche aufgetragen werden und entfaltet hier eine kühlende Wirkung. Die Pflanze wird auch bei Durchfall eingesetzt.
        \item \textbf{Bernstein} (20 Oth), \textit{Zutaten}.
        Ein glatter, ovaler Bernstein mit einem warmen, goldenen Schimmer. Seine polierte Oberfläche ist leicht transparent und reflektiert das Licht auf faszinierende Weise. Der handgroße Stein wirkt durch seine geschwungene Form wie ein natürlicher Talisman.
        \item \textbf{Kiefernharz} (10 Oth), \textit{Zutaten}.
        Das schützende Blut des Baumes. Wenn die raue Rinde einer Kiefer,sei es durch den Biss eines Wildtieres oder das ungeschickte Beil eines Holzfällers, verletzt wird, tritt diese zähe, goldgelbe Flüssigkeit aus.

Sie fließt langsam, beinahe träge, und erfüllt die Umgebung mit diesem unverwechselbaren, herben Duft.

Es wird genutzt, um spröde Korken zu versiegeln, Fackeln eine kräftige Flamme zu verleihen oder um einfache Wundsalben dickflüssig zu machen.

Sobald es an der Luft trocknet, wird es steinhart und schließt die Wunde des Baumes wie ein natürlicher Pfropfen. 

Wenn Ihr das Zeug erst einmal an den Händen habt, hilft nur noch ordentlich Fett oder sehr viel Geduld, um den Kleber wieder loszuwerden.
    \end{denseitemize}


    \subsection*{Trank der Macht}

    \begin{tabular}{@{}ll@{}}
        \textbf{\trans{difficulty_label}:} & Schwer \\
        \textbf{\trans{yield_label}:} & 1 \\
    \end{tabular}

    Zubereitung: Zuerst muss die Galle eines Bergwyverns in einem Tiegel aus gehärtetem Stahl langsam erhitzt werden, bis die ätzenden Dämpfe verflogen sind und eine goldgelbe Grundsubstanz zurückbleibt. In diese Basis werden die geraspelten Klauenstücke des Felsentrolls gegeben, welche die rohe, physische Kraft binden. Unter ständigem Rühren muss nun das getrocknete Tirakan-Moos hinzugefügt werden, welches als katalytischer Leiter dient. Sobald die Mischung zu leicht zu pulsieren beginnt, werden die Silberspäne eingestreut. Der Trank muss nun für genau eine Stunde über einem Feuer sieden, bis er dickflüssig wird und eine metallisch glänzende Farbe annimmt. Nach der Ruhezeit von einer Stunde kann er verwendet werden.

Zubereitungszeit: 3 Stunden (aktives Brauen), 12 Stunde Ruhezeit

Menge: 1 Anwendung, 200 Oth Flüssigkeit, 250 Oth Gewicht

Wirkung: Nach der Einnahme spürt der Charakter, wie sich seine Muskeln verhärten und seine Sinne schärfen. Die Anzahl aller Würfel die der Spieler wirft werden verdoppelt. Der Trank hält 2W6 Minuten an. Der Anwender ignoriert zudem alle Mali durch Erschöpfung oder Wunden und erhält einen temporären Rüstungsschutz von 2 gegen normalen Schaden (2xR) physischen Schaden (dieser addiert sich zu vorhandener Rüstung).

Sobald die Wirkung nachlässt, fordert die Magie ihren Tribut: Der Charakter erleidet sofort 2W6+6 Runden Erschöpfung und ist während diesem Zeitraum benommen (+2 Malus auf alle Mindestwürfe).

    \paragraph{\trans{ingredients_label}}
    \begin{denseitemize}
        \item \textbf{Bergwyvern Galle} (200 Oth), \textit{Zutaten}.
        Die Galle ist eine hochviskose, leuchtend goldgelbe bis giftgrüne Flüssigkeit, die in der Gallenblase der Bergwyvern gespeichert wird.

Der Geruch ist beißend, schwefelig und mit einer Note von verbranntem Kupfer. Schon das Einatmen der reinen Dämpfe kann die Nasenschleimhäute verätzen.

In Tiraka glaubt man, dass die Galle die „Essenz des ungestillten Hungers“ enthält. In der Alchemie wird sie als Katalysator verwendet, um andere Zutaten gewaltsam miteinander zu verschmelzen, die sich normalerweise abstoßen würden.

Die Gewinnung der Galle ist ein schwieriges Unterfangen für einen Alchemisten, da sie höchste Präzision unter widrigen Bedingungen erfordert.

Die Galle muss innerhalb von 1W6 Stunden nach dem Tod der Kreatur entnommen werden. Danach beginnt die Zersetzung der Gallenblase, und die Flüssigkeit verliert ihre alchemistische Potenz.

Man benötigt chirurgisches Besteck aus gehärtetem Stahl oder speziellen Keramikmessern. Einfaches Eisen würde von der Säure binnen Sekunden zerfressen werden.

Der Prozess des Eingriffs läuft wie folgt ab: Der Kadaver muss auf dem Rücken fixiert werden.
Ein tiefer Schnitt unterhalb des Brustbeins legt die Leber frei.
Die Gallenblase ist ein prall gefüllter, pulsierender Beutel. Sie muss mit einer Klemme am oberen Ende abgebunden werden, bevor man sie vorsichtig herausschneidet. Ein erfolgreicher Wurf auf Geschick (oder Medizin) gegen MW 8 ist erforderlich.

Sollte der Beutel platzen, erleidet der Erntende sofort 1W6 Schaden durch Verätzung, und die Zutat ist unwiederbringlich verloren. 

Wyvern-Galle kann nicht in normalen Glasfläschchen aufbewahrt werden, da sie das Glas mit der Zeit „blind“ macht und spröde werden lässt. Erfahrene Abenteurer nutzen bleigefütterte Tonkrüge oder Gefäße aus reinem Quarz, um die Substanz sicher nach Hause zu bringen.
        \item \textbf{Felsenmoos} (2 Bund), \textit{Zutaten}.
        Das Tirakan-Moos, von den Bergvölkern oft „Schattenflor“ oder „Schattensamt“ genannt, wächst meist in der Nähe von Bergwyvern Kollonien. Es ist in Felsspalten und kleinen höhlen zu finden. Man erzählt sich, dass es möglicherweise die Luft in engen Höhlen reinigt.

Das Moos wächst in dichten, schwammartigen Polstern. Seine Farbe ist ein tiefes, fast unnatürliches Dunkelviolett, das bei Berührung oder Luftzug in einem sanften, pulsierenden Indigo leuchtet (Biolumineszenz).
Bei Berührung hinterlässt es einen leicht klebrigen, metallisch riechenden Film auf der Haut.

Seine Hauptfunktion in Tränken ist die Erdung. Es verhindert, dass Energien oder rohe Kräfte sich verflüchtigen.

Roh zerkaut wirkt es stark schmerzlindernd und fiebersenkend, führt jedoch bei Überdosierung zu einer gefährlichen Verlangsamung des Herzschlags.

Die Ältesten behaupten, das Moos würde die Flüstern der Berge aufsaugen. Wer sein Ohr lange genug an ein Moospolster presst, soll die Stimmen der Vorfahren hören können oder den Berg, der vor Hunger knurrt.
        \item \textbf{Felsentroll-Klauen (Lithocrinus tirakanis)} (40 Oth), \textit{Zutaten}.
        Die Klauen eines Felsentrolls sind keine Krallen im biologischen Sinne, sondern mineralisierte Auswüchse aus verhärtetem Keratin und konzentrierten Erzen.

Sie spiegeln die unbändige, physische Regenerationskraft der Trolle wider, deren Haut eine Symbiose mit dem Gestein der Schattenfelsen eingegangen ist.

Eine einzelne, intakte Klaue hat etwa die Größe eines Kurzschwerters.
Die Färbung reicht von tiefem Grau bis hin zu Obsidian-Schwarz, wobei oft Schichten erkennbar sind, die an Schiefer erinnern.

Sie sind so hart, dass sie Funken schlagen, wenn sie auf Metall treffen. Gewöhnliche Klingen stumpfen an ihnen meistens sofort ab.

In der Alchemie wird fast nie die ganze Klaue verwendet, sondern eine verarbeitete Form. Die Klaue muss mühsam mit Diamantfeilen oder Meißeln aus gehärtetem Stahl bearbeitet werden. 

Als Zutat dienen fein geraspelte Späne oder Staub. Dieser Staub ist schwer, aschefarben und glitzert unter Lichteinfall.
Der Staub muss extrem fein sein. Zu grobe Späne lösen sich im Trank nicht auf und können beim Konsum die Speiseröhre des Anwenders schwer verletzen. Alternativ ist es ratsam, nach dem Brauvorgang, die Flüssigkeit durch ein Sieb zu geben.
        \item \textbf{Silberspähne} (10 Oth), \textit{Zutaten}.
        In Tiraka werden Silberspäne meist als Nebenprodukt in Schmieden oder bei der Herstellung von Schmuck gewonnen.

Für alchemistische Zwecke werden sie oft im Feuer gereinigt, um als Läutersilber die magischen Ströme in Tränken zu stabilisieren.
Silberspäne dienen als energetischer Anker. Sie verhindern, dass die instabilen Komponenten den Trank während des Brauprozesses zerreißen.
    \end{denseitemize}


    \subsection*{Meridischer Klarblick}

    \begin{tabular}{@{}ll@{}}
        \textbf{\trans{difficulty_label}:} & Mittel \\
        \textbf{\trans{yield_label}:} & 2 \\
    \end{tabular}

    Zubereitung:
Die Silberdistel muss zunächst in reinem Wasser aufgekocht und durch feinsten Kohlestaub gegossen werden, um die Wirkstoffe zu aktivieren und nur die zu erhalten die benötigt werden.

Anschließend werden die Mondbeeren vorsichtig untergerührt. Achtung: Die Mondbeeren dürfen nicht zerquetscht werden, Sie sollen nur ziehen und das 12 Stunden lang.

Zubereitungszeit:
4 Stunden aktive Destillation und Filterung an den Apparaturen, danach muss der Sud exakt 12 Stunden ruhen.

Menge:
2 Anwendungen. 400 Oth Flüssigkeit, Gesamtgewicht 600 Oth (inklusive dickwandiger Glasphiole).

Wirkung:
Neutralisiert Gifte im Blutkreislauf und schärft die Sinne, um den Geist vor Halluzinationen zu schützen. Bei Einnahme einer Anwendung wird der Status Vergiftet augenblicklich neutralisiert. Zusätzlich erhält der Konsument für die nächsten 1W6 Stunden einen Bonus von +1 Würfeln auf jegliche Resistenzproben gegen bewusstseinsverändernde Effekte oder Halluzinationen.

    \paragraph{\trans{ingredients_label}}
    \begin{denseitemize}
        \item \textbf{Kohle} (100 Oth), \textit{Zutaten}.
        Ein rein tiefschwarzer Stoff, der durch das kontrollierte Verkohlen von hartem Holz unter Luftabschluss gewonnen wird. Kohle.

Sie dient alchemistisch als universelles Filtriermedium. In Rezepte eingebunden oder zu feinstem Staub zerstoßen, saugt sie Verunreinigungen, Trübungen und unerwünschte Bitterstoffe wie ein Schwamm aus siedenden Flüssigkeiten.  

Sie ist überall in Tirakan bei Köhlern, Hufschmieden und gewöhnlichen Krämern in rauen Mengen für wenige Deut oder Gulden auf den städtischen Märkten erhältlich.
        \item \textbf{Silberdistel (Carlina acaulis)} (2 Bund), \textit{Zutaten}.
        Eine tief am Boden wachsende, widerstandsfähige Pflanze mit einer silbrig glänzenden, stacheligen Blütenkrone, die selbst bei eisigen Temperaturen nicht verkümmerte.

In der Alchemie gilt sie als reinigender Katalysator, dessen Essenz Gifte im Blut bindet, Entzündungen hemmt und die natürlichen Abwehrkräfte des Körpers stärkt.  

Sie gedeiht vor allem auf den kargen, sonnengefluteten Bergwiesen und felsigen Hängen des Toraner Hochlands. Man findet sie jedoch auch in gemässigteren Gebieten in ganz Tirakan.
        \item \textbf{Monddornbeere} (50 Oth), \textit{Zutaten}.
        Die Monddornbeere ist ein Geschenk der tiefsten Nacht. Sie wächst als bodendeckender Strauch, dessen zarte Ranken und tiefgrüne Blätter von markanten, kurzen Dornen geschützt werden.

Man findet sie ausschließlich an Orten, die selten das Licht der Sonne erblicken, meist tief in uralten Wäldern oder in der Nähe von feuchten Grotten- und Höhleneingängen. Ihre wahre Pracht entfaltet sie nur unter dem Licht des Vollmonds, wenn ihre kleinen, beerenartigen Früchte in einem geheimnisvollen, matten Blau leuchten, fast als hätten sie das Licht der Himmelskugel selbst verschluckt.

Die Beere ist berüchtigt für ihre starke sedierende Wirkung. In kleiner Dosis wirkt sie beruhigend, doch konzentriert in einem Trank, kann ihre Essenz den Geist betäuben und den Körper in einen Zustand tiefer, traumloser Stille versetzen. Das Sammeln ist riskant, da ihre Dornen bei Berührung einen temporären Juckreiz auslösen können.
        \item \textbf{Wasser} (350 Oth), \textit{Zutaten}.
        Kaltes klares Wasser.
    \end{denseitemize}


    \section{Gifte und Drogen}



    \subsection*{Linyas Tiefer Schlaf}

    \begin{tabular}{@{}ll@{}}
        \textbf{\trans{difficulty_label}:} & Schwer \\
        \textbf{\trans{yield_label}:} & 1 \\
    \end{tabular}

    Zubereitung:
Reibt den getrockneten Totenmohn zusammen mit den Baldrianwurzeln im Steinmörser zu einem homogenen, schwarzen Pulver. Setzt dieses in 150 Oth reinem Quellwasser an und lasst es zwei Stunden ungestört ziehen.

Siebt die Pflanzenreste durch ein feines Leinentuch in den kleinen Eisenkessel. Fügt die Melissenblätter zur Stabilisierung des organischen Äthers hinzu.

Erhitzt den Kessel (am besten auf dem Tisch des Brauers) unter stetigem Rühren, bis die Dämpfe bläulich schimmern. Fängt das Kondensat in einer Phiole auf.

Zubereitungszeit: 2 Stunden aktive Ruhezeit (Ziehenlassen), 3 Stunden aktives Sieden und Destillieren.

Menge: 1 Anwendung (ca. 50 Oth Flüssigkeit)

Wirkung: Wird die Essenz eingenommen oder über eine Waffe mit Giftkanal in die Blutbahn injiziert, verfällt das Opfer augenblicklich in den Zustand Bewusstlos.

Das Opfer muss zu Beginn jeder seiner Kampfrunden auf seine Willenskraft würfeln. Erst wenn die Anzahl der Erfolge (Würfe von 5+, explodierende 6en erlaubt) exakt 5 erreicht oder übertrifft, erwacht die Kreatur. Äußere Reize wie Rütteln oder Lärm sind wirkungslos; lediglich echter erlittener Schaden gewährt einen sofortigen Bonus von +2 Würfeln auf den Rettungswurf.

    \paragraph{\trans{ingredients_label}}
    \begin{denseitemize}
        \item \textbf{Totenmohn} (2 Bund), \textit{Zutaten}.
        Der Totenmohn wird ausschließlich auf der asgoraner Insel Linya angebaut, das Gewächs lässt sich sonst nirgends kultivieren. Bei der Pflanze handelt es sich um eine etwa zwanzig Finger hohe, mohnartige Blüte, welche zum Teil schwarz gefärbt ist. 

Der Totenmohn wird für allerhand Rituale und Tränke verwendet, was die Pflanze zu einem wahren Exportschlager Asgorans werden lässt.
        \item \textbf{Baldrian (Valeriana officinalis)} (1 Bund), \textit{Zutaten}.
        Baldrian hilft bei Einschlafstörungen und Unruhe. Hopfen und Melisse steigern die Wirkung des Baldrians und verbessern den Geschmack.
        \item \textbf{Melisse (Melissa officinalis)} (1 Bund), \textit{Zutaten}.
        Melisse wurde seit jeher als Heilkraut in der Medizin genutzt. Sie wirkt gegen Kopfschmerzen, Nervosität, Schlafstörungen und Magen-Darm-Beschwerden. Zudem bringt ein Aufguss mit Melisse Entspannung.
        \item \textbf{Wasser} (150 Oth), \textit{Zutaten}.
        Kaltes klares Wasser.
    \end{denseitemize}


    \subsection*{Morphin}

    \begin{tabular}{@{}ll@{}}
        \textbf{\trans{difficulty_label}:} & Schwer \\
        \textbf{\trans{yield_label}:} & 1 \\
    \end{tabular}

    Zubereitung:
Als natürliches Schmerzmittel findet sich Morphin primär im Schlafmohn (Papaver somniferum). Es ist der wichtigste Bestandteil des Opiums, das durch das Trocknen des Milchsafts aus den unreifen Samenkapseln gewonnen wird. Um den Wirkstoff zu isolieren, wird der wässrige Auszug mit Calciumchlorid oder Ammoniak behandelt. Dabei trennt sich das mekonsaure Calcium ab. Durch das spätere Eindampfen der verbleibenden Flüssigkeit gewinnt man schließlich Morphin und Codein in Form von Hydrochloridsalzen.

Zubereitungszeit: 60 Minuten zur Extraktion, danach 1 Stunde zum Abscheiden des Morphin.

Menge: 1 Anwendung, ca. 50 ml Flüssigkeit, Gewicht: 50 g.

Wirkung:
Morphin wirkt direkt im Gehirn und lindert dort das Schmerzempfinden sowie die Weiterleitung von Schmerzsignalen. Zusätzlich unterdrückt es den Hustenreiz und wirkt beruhigend auf den Körper, was auch den Blutdruck und den Puls senken kann.

Für Menschen, die das Medikament nicht gewohnt sind, können bereits 100 mg zu einer schweren Vergiftung führen. Wer das Mittel jedoch dauerhaft einnimmt, verträgt oft deutlich höhere Mengen. Morphin besitzt zudem ein hohes Suchtpotenzial. Ein häufiges Vorurteil ist, dass Morphin das Sterben beschleunigt. Dafür gibt es keine wissenschaftlichen Belege. Oft werden lediglich Begleiterscheinungen wie starke Schläfrigkeit oder Verwirrtheit fälschlicherweise als Zeichen für ein baldiges Lebensende gedeutet.

    \paragraph{\trans{ingredients_label}}
    \begin{denseitemize}
        \item \textbf{Calciumchlorid} (50 ml (Milliliter)), \textit{Zutaten}.
        Bei Calciumchlorid handelt es sich um eine Verbindung aus Calcium und Chlor mit der Formel CaCl\_2.

Chemisch gesehen befindet sich das Calcium darin in der Oxidationsstufe +2 und das Chlor in der Stufe -1. Man findet diesen Stoff in der Natur vor allem in Salzsolen gelöst. Wenn Calciumchlorid Wasser bindet und fest wird, entstehen daraus sehr seltene Minerale: Das Dihydrat wird als Sinjarit bezeichnet, während das Hexahydrat unter dem Namen Antarcticit bekannt ist. Man kann es jedoch auch einfach online bestellen.
        \item \textbf{Ammoniak} (10 ml (Milliliter)), \textit{Zutaten}.
        Ammoniak (NH\_3) fungiert als essenzieller Basischmikalie in der chemischen Industrie.

Das Hauptanwendungsgebiet liegt in der Produktion von stickstoffhaltigen Düngemitteln (z. B. Harnstoff oder Ammoniumnitrat), welche den Großteil des weltweiten Verbrauchs ausmachen. Darüber hinaus dient Ammoniak als Vorläufer für organische Synthesen, insbesondere bei der Herstellung von Polymeren und synthetischen Fasern.

Die Substanz ist als toxisch eingestuft. In hoher Konzentration wirken Ammoniakdämpfe stark reizend auf die Okularschleimhäute sowie das respiratorische System, wobei schwere Lungenschäden bis hin zum letalen Ausgang möglich sind. In der täglichen Anwendung sind solche Extremereignisse jedoch selten.
        \item \textbf{Schlafmohn (Papaver somniferum)} (5 Bund), \textit{Zutaten}.
        Der Schlafmohn (Papaver somniferum) gehört zur Familie der Mohngewächse und blickt auf eine lange Geschichte als eine unserer ältesten Heilpflanzen zurück.

Er ist jedoch nicht nur medizinisch interessant: Seine Samen werden in der Küche als Nahrungsmittel geschätzt oder zur Gewinnung von hochwertigem Speiseöl gepresst.

Das Besondere am Schlafmohn ist, dass fast alle Teile der Pflanze Alkaloide wie Morphin enthalten. Diese Wirkstoffe konzentrieren sich vor allem im weißen Milchsaft, der durch ein feines Kanalsystem die gesamte Pflanze durchläuft – besonders dicht ist dieses Netz in der Wand der Kapselfrucht. Um an diese Stoffe zu gelangen, werden die unreifen Kapseln leicht eingeritzt, sodass der Saft austreten kann. Wenn dieser Milchsaft an der Luft trocknet, entsteht das bekannte Betäubungsmittel Opium. Der Name leitet sich übrigens aus dem Griechischen ab und bedeutet schlicht „Säftchen“.
    \end{denseitemize}


    \subsection*{Flugschlangengift}

    \begin{tabular}{@{}ll@{}}
        \textbf{\trans{difficulty_label}:} & Einfach \\
    \end{tabular}

    Zubereitung:
Zuerst müsst Ihr das Rohgift der Flugschlange in einer kleinen Schale vorsichtig erwärmen. Aber hütet Euch vor Siedehitze! Wenn es zu dampfen beginnt, verfliegen die flüchtigen Lähmstoffe. Gebt das zerstoßene Schöllkraut hinzu, um die Durchblutung an der Einstichstelle beim Opfer zu reizen, so schlüpft das Gift schneller in die Säfte des Körpers. Rührt den getrockneten Salbei ein, bis die Flüssigkeit eine trübe, bernsteinartige Farbe annimmt. Zum Schluss wird das Ganze durch ein feines Leinentuch oder durch einen Filter aus Kohlenstaub gepresst, um die groben Rückstände zu entfernen. Übrig bleibt eine klare, leicht zähe Flüssigkeit.

Zubereitungszeit: 30 Minuten aktive Arbeit, danach etwa 1 Stunde Ruhezeit zum Abkühlen und Absetzen.

Menge: 1 Anwendung (ca. 5 Oth Flüssigkeit), Gesamtgewicht ca. 5 Oth.

Wirkung:
Das Opfer verspürt eine sofortige Kälte an der Wunde, die sich wie Blei in die Gliedmaßen ausbreitet. Es ist nicht tödlich, führt aber zu starker Benommenheit und verlangsamten Reflexen.
Das Opfer muss eine Probe auf Resistenz gegen MW 8 bestehen. Bei Misserfolg sinkt die Bewegungsweite für 1D6 Kampfrunden um 2 Punkte. Zudem erleidet das Opfer 1D3 Punkte Erschöpfung.

    \paragraph{\trans{ingredients_label}}
    \begin{denseitemize}
        \item \textbf{Salbei (Salvia officinalis)} (1 Bund), \textit{Zutaten}.
        Die Blätter des Salbeis haben eine entzündungshemmende, schweißhemmende und zusammenziehende Wirkung. Ein Tee oder Spülungen sind bei Halsentzündungen oder auch Schweißausbrüchen empfehlenswert.
        \item \textbf{Schöllkraut (Chelidonium majus)} (1 Bund), \textit{Zutaten}.
        Im Mittelalter wurde Schöllkraut bei Hautausschlägen, Sehschwäche oder Gelbsucht verwendet. Die Alkaloide der Pflanze habenden einen krampflösenden Effekt. Sie helfen bei Verdauungsproblemen und regen den Gallenfluss an.
        \item \textbf{Phiole Flugschlangengift} (1 Phiole), \textit{Tränke und Gifte}.
        Eine Phiole gefüllt mit Gift einer Flugschlange.
    \end{denseitemize}

\end{multicols}